% Options for packages loaded elsewhere
\PassOptionsToPackage{unicode}{hyperref}
\PassOptionsToPackage{hyphens}{url}
%
\documentclass[
]{article}
\usepackage{amsmath,amssymb}
\usepackage{iftex}
\ifPDFTeX
  \usepackage[T1]{fontenc}
  \usepackage[utf8]{inputenc}
  \usepackage{textcomp} % provide euro and other symbols
\else % if luatex or xetex
  \usepackage{unicode-math} % this also loads fontspec
  \defaultfontfeatures{Scale=MatchLowercase}
  \defaultfontfeatures[\rmfamily]{Ligatures=TeX,Scale=1}
\fi
\usepackage{lmodern}
\ifPDFTeX\else
  % xetex/luatex font selection
\fi
% Use upquote if available, for straight quotes in verbatim environments
\IfFileExists{upquote.sty}{\usepackage{upquote}}{}
\IfFileExists{microtype.sty}{% use microtype if available
  \usepackage[]{microtype}
  \UseMicrotypeSet[protrusion]{basicmath} % disable protrusion for tt fonts
}{}
\makeatletter
\@ifundefined{KOMAClassName}{% if non-KOMA class
  \IfFileExists{parskip.sty}{%
    \usepackage{parskip}
  }{% else
    \setlength{\parindent}{0pt}
    \setlength{\parskip}{6pt plus 2pt minus 1pt}}
}{% if KOMA class
  \KOMAoptions{parskip=half}}
\makeatother
\usepackage{xcolor}
\usepackage[margin=1in]{geometry}
\usepackage{graphicx}
\makeatletter
\def\maxwidth{\ifdim\Gin@nat@width>\linewidth\linewidth\else\Gin@nat@width\fi}
\def\maxheight{\ifdim\Gin@nat@height>\textheight\textheight\else\Gin@nat@height\fi}
\makeatother
% Scale images if necessary, so that they will not overflow the page
% margins by default, and it is still possible to overwrite the defaults
% using explicit options in \includegraphics[width, height, ...]{}
\setkeys{Gin}{width=\maxwidth,height=\maxheight,keepaspectratio}
% Set default figure placement to htbp
\makeatletter
\def\fps@figure{htbp}
\makeatother
\setlength{\emergencystretch}{3em} % prevent overfull lines
\providecommand{\tightlist}{%
  \setlength{\itemsep}{0pt}\setlength{\parskip}{0pt}}
\setcounter{secnumdepth}{-\maxdimen} % remove section numbering
\usepackage{amsmath}
\usepackage{mathtools}
\usepackage{amsthm}
\usepackage{multirow}
\usepackage{booktabs}
\ifLuaTeX
  \usepackage{selnolig}  % disable illegal ligatures
\fi
\IfFileExists{bookmark.sty}{\usepackage{bookmark}}{\usepackage{hyperref}}
\IfFileExists{xurl.sty}{\usepackage{xurl}}{} % add URL line breaks if available
\urlstyle{same}
\hypersetup{
  pdftitle={Discrete Multivariate Analysis - 579 - HW \#7},
  pdfauthor={Alex Towell (atowell@siue.edu)},
  hidelinks,
  pdfcreator={LaTeX via pandoc}}

\title{Discrete Multivariate Analysis - 579 - HW \#7}
\author{Alex Towell
(\href{mailto:atowell@siue.edu}{\nolinkurl{atowell@siue.edu}})}
\date{}

\begin{document}
\maketitle

\newcommand{\var}{\operatorname{Var}}
\newcommand{\expect}{\operatorname{E}}
\newcommand{\corr}{\operatorname{Corr}}
\newcommand{\cov}{\operatorname{Cov}}
\newcommand{\ssr}{\operatorname{SSR}}
\newcommand{\se}{\operatorname{SE}}
\newcommand{\mat}[1]{\mathbf{#1}}
\newcommand{\RR}{\operatorname{RR}}
\newcommand{\eval}[2]{\left. #1 \right\vert_{#2}}

\hypertarget{problem-1}{%
\section{Problem 1}\label{problem-1}}

Consider an experiment on chlorophyll inheritance in maize. A genetic
theory predicts the ratio of green to yellow to be 3:1. In a sample of
\(n = 1103\) seedlings, \(n_1 = 854\) were green and \(n_2 = 249\) were
yellow.

\hypertarget{part-a}{%
\subsection{Part (a)}\label{part-a}}

\fbox{Compute the statistic $X^2$ for testing the proposed model.}

The statistic is defined as \[
  X^2 = \sum_{j=1}^{c} \frac{(n_j - n \pi_{j 0})^2}{n \pi_{j 0}}.
\]

Under the null model, the odds are 3:1, or \(\pi_{1 0} = 0.75\) and
\(\pi_{1 1} = 0.25\). We are given \(n = 1103\), \(n_1 = 854\), and
\(n_2 = 249\). Under the null model, these values have expectations
given respectively by \(n \pi_{1 0} = 827.25\) and \(275.75\).

The observed statistic is thus given by \[
  X_0^2 = \frac{(854 - 827.25)^2}{827.25} + \frac{(249 - 275.75)^2}{275.75} = 3.46.
\]

\hypertarget{part-b}{%
\subsection{Part (b)}\label{part-b}}

\fbox{Determine the upper 10th percentile for the reference distribution.}

The reference distribution is the chi-squared distribution with \(1\)
degree of freedom, denoted by \(\chi^2(1)\).

The upper 10th percentile given \(1\) degree of freedom, denoted by
\(\chi_{0.10}^2(1)\), is found by solving for \(\chi_{0.10}^2(1)\) in
the equation \(\Pr(\chi^2(1) \geq \chi_{0.10}^2(1)) = 1-0.10 = 0.9\),
which yields the result \[
  \chi_{0.10}^2(1) = 2.71.
\]

\hypertarget{part-c}{%
\subsection{Part (c)}\label{part-c}}

\fbox{Provide an interpretation of your result, stated in the context of the problem.}

We see that any observed statistic \(X_0^2\) with \(\rm{df}=1\) greater
than \(\chi^2_{0.10}(1) = 2.71\) is not compatible with the null model
at significance level \(\alpha = 0.10\).

Since the observed statistic \(X_0^2 = 3.46 > 2.71\), the null model,
which is the genetic theory where the ratio of green to yellow is 3:1,
is not compatible with the data.

\hypertarget{part-d}{%
\subsection{Part (d)}\label{part-d}}

\fbox{What are the shortcomings of hypothesis testing as a measure of evidence?}

Hypothesis testing, as a dichotomous measure of evidence, does not
provide as much information as a more quantitative evidence measure. For
instance, it does not provide information about \emph{effect size}.

\hypertarget{problem-2}{%
\section{Problem 2}\label{problem-2}}

\hypertarget{part-a-1}{%
\subsection{Part (a)}\label{part-a-1}}

\fbox{Compute a $90\%$ confidence interval for $\pi_1$.}

The MLE of \(\pi_1\) is given by
\(\hat{\pi}_1 = \frac{n_1}{n} = \frac{854}{1103} = 0.774\). Letting
\(\alpha = 0.10\) and inverting the Wald test statistic, we get the
\(90\%\) confidence interval for \(\pi_1\), \[
  \hat{\pi}_1 \pm z_{1-alpha/2} \sigma_{\hat\pi} = 0.744 \pm 1.645 \sqrt{\frac{0.774(1-0.774)}{1103}},
\] which may be rewritten as \[
  [0.754, 0.795].
\]

\hypertarget{part-b-1}{%
\subsection{Part (b)}\label{part-b-1}}

\fbox{Provide an interpretation of your result, stated in the context of the problem.}

Based on the observed data, we estimate that the probability \(\pi_1\)
of a green strain is between \(0.754\) and \(0.795\).

As expected, the null model specifies a value for \(\pi_1\) (\(0.75\))
that is not contained in this confidence interval.

\end{document}
